\begin{pastexam}[label={px:ntu-109-comm-dsa-02}]{臺大電信所 109 學年度（B）第 2 題}
（是非；每題 2 分，答錯倒扣 1 分）

Depth-first search will take $\Theta(V^2)$ time on a graph $G=(V,E)$ represented as an adjacency
matrix.
\end{pastexam}

\begin{pastsoln}
答案：T（對）。

DFS 對每個頂點 $u$ 只呼叫一次 \proc{DFS-Visit}$(u)$，呼叫時要找出 $u$ 的所有鄰居。 用鄰接
矩陣存，找鄰居就是把第 $u$ 列的 $|V|$ 格全部看一遍，不管這一列有幾個 $1$。 所以：
\begin{itemize}
\item 初始化（所有頂點設成白色）：$\Theta(V)$。
\item 每個頂點掃一整列：$|V|$ 個頂點 $\times$ $|V|$ 格 $= \Theta(V^2)$。
\end{itemize}
合計 $\Theta(V^2)$，而且上下界都是 $V^2$：就算圖裡一條邊都沒有，每一列還是要看完才知道
沒有鄰居。 用鄰接串列才是 $\Theta(V + E)$（資結書第~21.2~節的比較表）。
\end{pastsoln}
